% Part: lambda-calculus
% Chapter: syntax
% Section: abbreviated-syntax

\documentclass[../../../include/open-logic-section]{subfiles}

\begin{document}

\olfileid{lam}{syn}{abb}
\olsection{Sintaksis kang Dicekak}

Suku kang ditetepake ing
\olref[trm]{defn:term} kadhangkala
ruwet yen ditulis, mula migunani
yen kita ngenalake sintaksis kang
luwih ringkes. Mesthi wae kita kudu
njamin yen suku kang ditulis kanthi
notasi ringkes iku tetep bisa diwaca
kanthi cara tunggal. Salah siji
versi kang lumrah digunakake,
diarani \emph{suku cekakan},
kaya mangkene.

\begin{enumerate}
\item Yen tandha kurung ora ditulis,
  aplikasi lumaku saka kiwa menyang
  tengen. Umpamane, yen $M$, $N$,
  $P$, lan~$Q$ iku suku, mula
  $MNPQ$ iku cekakan saka
  $(((MN)P)Q)$.
\item Maneh, yen tandha kurung ora
  ditulis, abstraksi lambda nduweni
  cakupan paling amba kang bisa.
  Umpamane, $\lambd[x][MNP]$
  diwaca minangka
  $(\lambd[x][MNP])$.
\item Lambda bisa digunakake kanggo
  ngiket luwih saka siji variabel.
  Umpamane, $\lambd[xyz][M]$
  iku cekakan saka
  $\lambd[x][\lambd[y][\lambd[z][M]]]$.
\end{enumerate}

Umpamane,
\[
\lambd[xy][xxyx \lambd[z][xz]]
\]
iku cekakan saka
\[
(\lambd[x][(\lambd[y][((((xx)y)x)(\lambd[z][(xz)]))])]).
\]

\begin{prob}
Uraikna suku cekakan
$\lambd[g][(\lambd[x][g (x x)])
  \lambd[x][g (x x)]]$.
\end{prob}

\end{document}
